% English translation of questions/5784/semA-special/q4.tex (metadata is taken from the Hebrew file).

\begin{question}
Prove that the following equation has a unique positive real solution:
\[ 2x^3+\arctan x+\sqrt{x}=7 \]
\end{question}

\begin{hint}
Existence: the Intermediate Value Theorem for $g(x)=2x^3+\arctan x+\sqrt x-7$ on a suitable interval. Uniqueness: show that $g$ is strictly increasing.
\end{hint}

\begin{solution}
Define $g(x)=2x^3+\arctan x+\sqrt x-7$ on $[0,\infty)$. $g$ is continuous there as a sum of continuous functions, and differentiable on $(0,\infty)$.

\textbf{Existence.} $g(0)=-7<0$, and
\[ g(2)=16+\arctan2+\sqrt2-7=9+\arctan 2+\sqrt2>0. \]
By the Intermediate Value Theorem there exists $x_0\in(0,2)$ with $g(x_0)=0$, i.e. a positive solution of the equation.

\textbf{Uniqueness.} For every $x>0$:
\[ g'(x)=6x^2+\frac1{1+x^2}+\frac1{2\sqrt x}>0. \]
Hence $g$ is strictly increasing on $(0,\infty)$ (a corollary of Lagrange's theorem: if $0<x_1<x_2$ then $g(x_2)-g(x_1)=g'(c)(x_2-x_1)>0$). A strictly increasing function attains each value at most once, hence it has at most one zero in $(0,\infty)$.

(Alternatively: if there were two positive solutions $x_1<x_2$, then by Rolle's theorem there would be $c\in(x_1,x_2)$ with $g'(c)=0$, contradicting $g'>0$.)

\textbf{Conclusion:} the equation has a unique positive real solution (numerically $x_0\approx1.3486$).
\end{solution}
